\begin{proposition}[The global partial filling]\hypertarget{tag-030N}{}\label{push:src:B.4.19}
\label{push:B419}
\StatementRecord{push:src:B.4.19}{proved}{030N}{\SeriesRef{push:src:B.1.2}, \SeriesRef{push:src:B.4.18}, \SeriesRef{push:src:B.4.30}, \SeriesRef{push:src:B.4.33}, \SeriesRef{push:src:C.3.5}, \SeriesRef{push:src:B.1.1}, \SeriesRef{push:src:B.3.2}, \SeriesRef{push:src:B.4.9}, \SeriesRef{push:src:A.1.2}}

Let
\[
 \Omega=(\mathbb P^1\times\Delta_r)
                  \setminus\{(\infty,0)\}.
\]
There is a proper Hausdorff holomorphic map
$\mathcal Y^\circ\to\Omega$ with smooth total fourfold,
flat over $r$, whose restriction to $r\ne0$ is the
family pulled back by $w=r^6$. Its reduced
normal-crossings central divisor has the finite-place
models of \ref{push:src:B.1.2} and the punctured cusp of
\ref{push:src:B.4.18}, with their given markings. Over
$B^\circ=\mathbb P^1\setminus\{0,1,\infty\}$ its
relative component cover is
\[
 y^{12}=K(x),\qquad K(x)=\frac{x^4}{(x-1)^3}.
\]
Thus the global finite-place overlap maps and their
cocycle are constructed. The sole missing base point is $(\infty,0)$.
Proposition~\ref{push:src:B.4.30} proves its noncompletion, and
\ref{push:src:B.4.33} excludes every proper filling of the 
family over the whole $r$ disc.
\end{proposition}



\begin{proof}
\begin{enumerate}[label=\arabic*.,leftmargin=*]
\item We distinguish the ordinary base coordinate $x$ from
the logarithmic parameter
\[
 v=\frac{\log r}{2\pi i},\qquad c=-12v,
 \qquad e=(0,1)^t.
\]
Use the coherent closing of \ref{push:src:C.3.5}. Its additive formula
for a positive $r$ loop, that is $c\mapsto c-12$, is
\begin{equation}
 b_{-12,c}=\Pi_c\widehat\gamma-(k+6)e,
 \qquad k=12\ell=\frac{\log K}{2\pi i}.
\label{pub:I:eq:B-4-19-1}
\end{equation}
The first summand and $6e$ are periods.
Consequently its torus translation has the scalar
representative $-ke$, since addition of those periods
preserves the quotient map.

\item The coherent closing in \ref{push:src:C.3.5} may be used for the
specified punctured family. More explicitly, two
coherent choices differing by a scalar automorphism
give a principal $\mathbb C^*$ torsor over the punctured
parameter disc. That torsor is holomorphically trivial:
the exponential sequence and Cartan's theorem B give
$\operatorname{Pic}(\Delta_r^*)=H^2(\Delta_r^*,\mathbb Z)=0$.
The resulting scalar gauge is a holomorphic
isomorphism of the punctured families.

\item \emph{The ordinary charts and their exact cocycle.}
Choose a good analytic cover of $B^\circ$, refining
period-frame charts, and on each member choose a
holomorphic logarithm $k_i=(2\pi i)^{-1}\log K$.
On the parameter logarithm cover put
\begin{equation}
 \widetilde\zeta_i=\zeta+k_i v e.
\label{pub:I:eq:B-4-19-2}
\end{equation}
Under the positive $r$ deck, $\zeta$ changes by
$-k_i e$ modulo periods and $v$ changes by one.
Thus \eqref{pub:I:eq:B-4-19-2} has the ordinary period descent.
It identifies the local punctured family
with the standard twelve-cycle plumbing of \ref{push:src:B.1.1},
now over an ordinary base chart.

\item On an overlap let $n_{ij}=k_j-k_i\in\mathbb Z$.
The invariant vector $e$ is fixed by every period-frame
change. In the multiplicative coordinate
$W_i=\exp(-2\pi i\widetilde\zeta_{i,2})$, the extra
transition is exactly
\begin{equation}
 W_j=r^{-n_{ij}}W_i.
\label{pub:I:eq:B-4-19-3}
\end{equation}
Every integral power in \eqref{pub:I:eq:B-4-19-3} extends to an
automorphism of the whole twelve-cycle plumbing.
On the unfolded node chart write
\[
 x_j=W/r^j,\qquad y_j=r^{j+1}/W,
 \qquad x_jy_j=r.
\]
Multiplication $W\mapsto r^{-n}W$ sends the chart
index $j$ to $j-n$ and leaves both displayed chart
coordinates unchanged. The same statement holds
in bundle frames, whose additional factors are
holomorphic units. It commutes with the closing
period, which shifts the index by twelve and
translates the elliptic base by $6P$.
The usual semiabelian translations and period-frame
changes preserve the two sections, hence also extend.

\item The integers satisfy
$n_{ik}=n_{ij}+n_{jk}$. Since the scalar translations
commute with each other and with the period-frame
changes, \eqref{pub:I:eq:B-4-19-3} is an exact cocycle on the entire
plumbing. This constructs the family over
$B^\circ\times\Delta_r$, with the original
punctured-family identification \eqref{pub:I:eq:B-4-19-2}.

\item \emph{The finite-place cancellations.}
At the two root charts put
\[
 w_j=\frac{\log s_j}{2\pi i},\qquad
 \epsilon_1=1,\quad\epsilon_2=-1,
 \qquad k_j=12\kappa_j.
\]
The notation is that of \ref{push:src:C.3.5}: $k_j$ is a
holomorphic logarithm of the unit
$K(x)s_j^{-12\epsilon_j}$ on the full root disc.
Its invariance under the finite deck is exact.
Indeed the unit is invariant since $m_j\mid12$;
the difference of two logarithms is an integer
constant, and evaluating at the fixed root centre
makes that constant zero. The scalar representative
of the local positive $r$ closing is $-k_j e$.
Apply the local gauge
$\widetilde\zeta_j=\zeta_j+k_jve$ and use the
smooth twelve-cycle root model of \ref{push:src:B.1.1}, with its
free finite deck action from \ref{push:src:B.1.2}. Because $e$ is
invariant and $k_j$ is invariant, this gauge does
not change that finite action.

\item On a sector of the punctured root disc choose
compatible logarithms. The root-to-ordinary
collar is
\[
 \zeta=\zeta_j+\sigma_{j,c},\qquad
 \sigma_{j,c}=w_j\Pi_c v_j,
 \qquad \gamma(v_j)=\epsilon_j.
\]
Moreover $k-k_j=12\epsilon_jw_j$. Therefore its
expression in the new gauges is
\begin{equation}
 \begin{aligned}
 \widetilde\zeta-\widetilde\zeta_j
 &=w_j\Pi_c v_j+12\epsilon_jw_jve\\
 &=w_j\Pi_{\rm part}v_j,
 \qquad \Pi_{\rm part}=\Pi_c-ce\gamma.
 \end{aligned}
\label{pub:I:eq:B-4-19-4}
\end{equation}
Here $c=-12v$ is essential: the parameter logarithms
cancel exactly. The last translation is independent
of $r$. On each sector it is a holomorphic section
of the semiabelian group and therefore extends to
the ruled normalization and all node charts.

\item Changing a logarithm adds an integral transition
of the form \eqref{pub:I:eq:B-4-19-3}, together with a
period-frame identification. Those maps already
extend. Their deck and overlap identities agree
on the dense locus $r\ne0$ with the collar
identities, hence agree everywhere by the identity
principle. Thus the sector maps descend and glue
exactly to the finite quotients. This constructs
the specified family over
$(\mathbb P^1\setminus\{\infty\})\times\Delta_r$.
It retains the multiplicities three and four and
the reductions already proved in \ref{push:src:B.1.2}.

\item The component cover is determined by the local indices.
If $i_i$ is the
unfolded component index modulo twelve, then
\eqref{pub:I:eq:B-4-19-3} gives $i_j=i_i-n_{ij}$. Hence
\[
 y=\exp(2\pi i k_i/12)\exp(2\pi i i_i/12)
\]
is independent of the chart and satisfies $y^{12}=K$.
The clockwise deck shifts are consequently $4,-3,-1$.
The local finite quotients complete this cover at
$0,1$ exactly as in \ref{push:src:B.3.2}.

\item \emph{The cusp overlap, including $r=0$.}
Retain $s=\log q/(2\pi i)$ and the holomorphic
functions $C,\mu,\kappa_0$ of \ref{push:src:B.4.18} and \ref{push:src:C.3.5}.
Then $k=-s+12\kappa_0$. The cusp-to-rotating comparison of the positive-deck model is
$\zeta\mapsto\zeta-\sigma+F$. Substitution in the
displayed formulas of \ref{push:src:B.4.9} and \ref{push:src:B.4.18} gives
\begin{equation}
 F-\sigma=kve+G(s)+6ve,
 \qquad
 G(s)=\left(\frac{\mu s}{2},\,
                    \frac{Cs}{12}-\frac{s^2}{24}\right).
\label{pub:I:eq:B-4-19-5}
\end{equation}
Thus its comparison with the ordinary gauge
\eqref{pub:I:eq:B-4-19-2} is the translation $G(s)+6ve$.
On a sector of $q\ne0$, $G$ is independent of $r$
and gives a holomorphic semiabelian translation.
The other term is the integral fan shift
$W\mapsto r^{-6}W$. Both extend, with their inverses,
over the full twelve-cycle $r$-central divisor.
The branch identities are
\[
 G(s+1)-G(s)=b(s,v)+ve,
 \qquad k(s+1)-k(s)=-1.
\]
The ordinary coordinate deck contributes $-ve$;
their sum is exactly the positive fractional deck
$b$. A positive $r$ loop changes $6ve$ by the
integral period $6e$. These identities prove
descent of the extended comparison to the full
punctured $q$ disc. In particular the common
clockwise cusp shift is $-1$, as required.

\item Glue the just constructed global model on
$x\ne\infty$ to \ref{push:src:B.4.18} on the cusp neighbourhood
minus its origin, using \eqref{pub:I:eq:B-4-19-5}. The overlap
contains $r=0$ and the map extends there by the
preceding chart calculation. The identities give the exact scalar and translation cocycles
with the prescribed punctured marking.

\item Each chart used is a proper model over its base
open set: this was proved by compact-core plumbing
in \ref{push:src:A.1.2} and by the finite free quotient in \ref{push:src:B.1.2};
\ref{push:src:B.4.18} is proper on its stated punctured bidisc.
Properness is local on the target, so the glued
map to $\Omega$ is proper. The same local argument
gives Hausdorffness, smoothness of the total space,
flatness over $r$, and the reduced normal-crossings
central divisor. For $r\ne0$ the base is the entire
$\mathbb P^1$, and the specified comparison is
precisely the compact family. At $r=0$ the target is the finite base $\mathbb P^1\setminus\{\infty\}$,
as specified in $\Omega$.

\end{enumerate}
\end{proof}